% ============================================================
%  Dokumentacja Analityczno-Projektowa
%  Mobilny System Wspomagania Honorowych Dawców Krwi
% ============================================================
\documentclass[11pt, a4paper]{article}

% ── Encoding & Language ──────────────────────────────────────
\usepackage[T1]{fontenc}
\usepackage[utf8]{inputenc}
\usepackage[provide=*,polish]{babel}

% ── Typography ───────────────────────────────────────────────
\usepackage{palatino}
\usepackage{microtype}
\usepackage{setspace}
\setstretch{1.15}

% ── Page Layout ──────────────────────────────────────────────
\usepackage[
  top=25mm, bottom=28mm,
  left=28mm, right=28mm,
  headheight=14pt
]{geometry}

% ── Headers & Footers ────────────────────────────────────────
\usepackage{fancyhdr}
\pagestyle{fancy}
\fancyhf{}
\renewcommand{\headrulewidth}{0.3pt}
\fancyhead[L]{\small\textit{Dokumentacja Analityczno-Projektowa}}
\fancyhead[R]{\small\textit{System Krwiodawstwa }}
\fancyfoot[C]{\small\thepage}

% ── Section Styling ──────────────────────────────────────────
\usepackage{titlesec}
\titleformat{\section}
  {\normalfont\large\bfseries\scshape}
  {\thesection.}{0.6em}{}[\titlerule]
\titlespacing{\section}{0pt}{18pt}{8pt}

\titleformat{\subsection}
  {\normalfont\normalsize\bfseries}
  {\thesubsection.}{0.5em}{}
\titlespacing{\subsection}{0pt}{12pt}{4pt}

\titleformat{\subsubsection}
  {\normalfont\small\bfseries\itshape}
  {\thesubsubsection.}{0.5em}{}
\titlespacing{\subsubsection}{0pt}{8pt}{3pt}

% ── Tables ───────────────────────────────────────────────────
\usepackage{booktabs}
\usepackage{tabularx}
\usepackage{array}
\renewcommand{\arraystretch}{1.35}
\newcolumntype{L}[1]{>{\raggedright\arraybackslash}p{#1}}

% ── Requirement Cards ────────────────────────────────────────
\usepackage{mdframed}
\usepackage{xparse}

\mdfdefinestyle{reqcard}{
  linewidth=0.5pt,
  linecolor=black,
  backgroundcolor=white,
  innertopmargin=10pt,
  innerbottommargin=10pt,
  innerleftmargin=12pt,
  innerrightmargin=12pt,
  skipabove=10pt,
  skipbelow=6pt,
  roundcorner=2pt,
}

\mdfdefinestyle{reqheader}{
  linewidth=0.5pt,
  linecolor=black,
  backgroundcolor=black!8,
  innertopmargin=6pt,
  innerbottommargin=6pt,
  innerleftmargin=12pt,
  innerrightmargin=12pt,
  skipabove=0pt,
  skipbelow=0pt,
  roundcorner=2pt,
}

% ── Lists ────────────────────────────────────────────────────
\usepackage{enumitem}
\setlist[itemize]{leftmargin=1.4em, topsep=3pt, itemsep=2pt, parsep=0pt}
\setlist[enumerate]{leftmargin=1.6em, topsep=3pt, itemsep=2pt, parsep=0pt}

% ── Misc ─────────────────────────────────────────────────────
\usepackage{parskip}
\setlength{\parskip}{5pt}
\usepackage{graphicx}
\usepackage{xspace}
\usepackage[gray]{xcolor}
\usepackage{hyperref}
\hypersetup{colorlinks=false, pdfborder={0 0 0}}
\usepackage{caption}

% ── Helpers ──────────────────────────────────────────────────
\newcommand{\req}[1]{\texttt{#1}}
\newcommand{\term}[1]{\textbf{#1}}
\newcommand{\sys}{\textsc{hdk\,app}\xspace}

% ════════════════════════════════════════════════════════════
\begin{document}
% ════════════════════════════════════════════════════════════

% ── Title Page ───────────────────────────────────────────────
\begin{titlepage}
  \centering
  \vspace*{30mm}

  {\fontsize{9}{11}\selectfont\scshape dokumentacja analityczno-projektowa}\\[4mm]
  \rule{0.55\linewidth}{0.4pt}\\[6mm]

  {\fontsize{22}{28}\selectfont\bfseries
    Mobilny System Wspomagania\\[3mm]
    i Monitorowania Aktywności\\[3mm]
    Honorowych Dawców Krwi}\\[8mm]

  \rule{0.55\linewidth}{0.4pt}\\[6mm]

  {\large\itshape 
  Daniel Abbas 281134 \\[1mm]
                    Jakub Adamczyk 281206 \\[1mm]
                    Jakub Kret 281182\\[6mm]


  \begin{tabular}{rl}
    \toprule
    \textbf{Status} & Projekt wstępny \\
    \textbf{Data}   & 18.04.2026 \\
    \textbf{Platforma} & Android (Flutter) \\
    \bottomrule
  \end{tabular}

  \vfill
  {\small\scshape projekt zespołowy \,·\, informatyka techniczna}
\end{titlepage}

\newcommand{\letterspacingtext}[1]{#1}

% ── Table of Contents ────────────────────────────────────────
\tableofcontents
\thispagestyle{empty}
\newpage

% ════════════════════════════════════════════════════════════
\section{Wstęp i założenia organizacyjne projektu}
% ════════════════════════════════════════════════════════════

\subsection{Kontekst inicjacji i cele zespołu}

Głównym celem systemu jest stworzenie cyfrowego ekosystemu, który nie tylko ułatwia
logistykę oddawania krwi (śledzenie terminów, lokalizacja punktów pobrań), ale przede
wszystkim buduje trwałą relację z dawcą i promuje postawy prospołeczne.

Analiza dostępnych rozwiązań mobilnych wykazała, że brakuje na rynku narzędzi, które
realnie angażowałyby dawców i spełniały ich oczekiwania.
Aplikacja ma wypełnić tę lukę poprzez:

\begin{itemize}
  \item \textbf{Personalizację:} dostosowanie komunikatów do konkretnego użytkownika.
  \item \textbf{Grywalizację:} wprowadzenie mechanizmów zwiększających motywację i retencję dawców.
  \item \textbf{Użyteczność:} nowoczesne podejście do doświadczeń medycznych (UX w zdrowiu).
\end{itemize}

\subsubsection*{Metodyka i standardy}

\begin{itemize}
  \item \textbf{Zarządzanie:} uwzględniliśmy podział ról w zespole, analizę ryzyka oraz harmonogram prac.
  \item \textbf{Estymacja kosztów:} koszty produkcji oszacowano empirycznie.
  \item \textbf{Specyfikacja wymagań:} wymagania funkcjonalne i niefunkcjonalne opisaliśmy zgodnie
        z uznanym wzorcem \textbf{Volere}.
\end{itemize}

% ────────────────────────────────────────────────────────────
\subsection{Dekompozycja ról i zakresy odpowiedzialności w zespole}

\begin{center}
\begin{tabularx}{\linewidth}{L{3.8cm} X}
  \toprule
  \textbf{Rola w~Zespole} & \textbf{Zakres Odpowiedzialności i Zadania Techniczne} \\
  \midrule
  Lider Projektu / Architekt Frontendu i~UI/UX
    & Reprezentacja zespołu na zewnątrz, prowadzenie spotkań weekly, przesyłanie
      dokumentacji. Technicznie: projektowanie użyteczności aplikacji (UI/UX),
      implementacja widoków (Flutter), integracja liczników i map oraz projektowanie
      mechanizmów grywalizacyjnych. \\
  \addlinespace
  Główny Inżynier Backendowy / Architekt Algorytmów
    & Projektowanie i~zapewnienie spójności między logiką serwerową a~frontendem
      oraz skalowalność i~powtarzalność środowisk dzięki izolacji kontenerowej.
      Implementacja algorytmów obliczających karencję oraz kalkulatora ulgi podatkowej PIT.
      Badanie i~implementacja optycznego rozpoznawania znaków (OCR) dla skanera wyników
      badań morfologicznych. \\
  \addlinespace
  Inżynier Danych / Analityk Wymagań (Volere)
    & Projektowanie autoryzacji oraz architektury powiadomień Push.
      Odpowiedzialność za inżynierię wymagań metodą Volere.
      Administracja bazami danych. \\
  \bottomrule
\end{tabularx}
\end{center}

% ════════════════════════════════════════════════════════════
\section{Uzasadnienie biznesowe, cele i analiza tła dziedzinowego}
% ════════════════════════════════════════════════════════════

\subsection{Działalność użytkownika i kontekst problemu biznesowego}

Krew i jej składniki (osocze, płytki krwi) są niemożliwymi do sztucznego wyprodukowania,
krytycznymi zasobami terapeutycznymi. System krwiodawstwa w Polsce opiera się na idei
honorowego dawstwa, jednakże struktura demograficzna dawców ulega starzeniu, a pozyskiwanie
i — co ważniejsze — utrzymanie przedstawicieli młodszych pokoleń napotyka na znaczne
trudności.

Główne problemy systemowe identyfikowane z perspektywy dawcy:

\begin{enumerate}
  \item \textbf{Brak proaktywnej komunikacji:} Centra krwiodawstwa (RCKiK) rzadko informują
        o spersonalizowanym zapotrzebowaniu na konkretną grupę krwi w sposób natychmiastowy
        i angażujący.

  \item \textbf{Dezorientacja w zakresie uprawnień:} Przepisy dotyczące odstępów między
        donacjami są wysoce skomplikowane — różnią się w zależności od oddanego składnika
        (krew pełna: 8 tygodni, osocze: 2 lub 4 tygodnie, płytki krwi: 4 tygodnie) oraz
        płci (mężczyźni: do 6 donacji krwi pełnej rocznie, kobiety: do 4 donacji).
        Błędne wyliczenia prowadzą do wizyt zakończonych odrzuceniem medycznym.

  \item \textbf{Niski poziom cyfryzacji świadczeń:} Mimo postępów w usługach e-zdrowia,
        dawcy często muszą samodzielnie archiwizować papierowe wyniki badań, a rozliczenie
        ulgi podatkowej z tytułu darowizny krwi w PIT (130 zł za każdy oddany litr) wymaga
        manualnego prowadzenia skomplikowanych rejestrów.

  \item \textbf{Brak mechanizmów aktywnej motywacji zewnętrznej:} Badania naukowe wskazują,
        że brak elementów gamifikacji i natychmiastowych gratyfikacji psychologicznych
        powoduje spadek partycypacji po pierwszych donacjach.
\end{enumerate}

% ────────────────────────────────────────────────────────────
\subsection{Cele projektu}

\begin{itemize}
  \item \textbf{Zwiększenie częstotliwości i regularności donacji} o prognozowane 20\%
        w grupie docelowej poprzez implementację zaawansowanych liczników karencji,
        powiadomień celowanych (o brakach konkretnej grupy krwi) oraz mechanizmów
        śledzenia programu lojalnościowego.

  \item \textbf{Automatyzacja procesów logistycznych i administracyjnych} dla dawcy,
        w tym geolokalizacja stałych centrów RCKiK i mobilnych akcji poboru (krwiobusów)
        na zintegrowanej mapie oraz dostarczenie automatycznego kalkulatora ulgi
        podatkowej PIT.

  \item \textbf{Zbudowanie spersonalizowanego profilu zdrowotnego,} chroniącego dawcę
        przed przeciwwskazaniami fizjologicznymi. Cel ten zostanie osiągnięty przez
        wdrożenie mechaniki monitorowania cyklu miesiączkowego u kobiet oraz cyfrowego
        archiwum badań z wykorzystaniem optycznego rozpoznawania znaków (OCR).
\end{itemize}

% ────────────────────────────────────────────────────────────
\subsection{Słownik pojęć, konwencje i regulacje dziedzinowe}

\begin{center}
\begin{tabularx}{\linewidth}{L{3cm} X}
  \toprule
  \textbf{Termin} & \textbf{Definicja} \\
  \midrule
  \term{HDK}
    & Honorowy Dawca Krwi — zarejestrowana w publicznej służbie krwi osoba
      oddająca bezpłatnie krew lub jej składnik. \\
  \addlinespace
  \term{ZHDK}
    & Zasłużony Honorowy Dawca Krwi — tytuł nadawany przez PCK.
      Progi litrażowe: kobiety 5\,L / 10\,L / 15\,L; mężczyźni 6\,L / 12\,L / 18\,L.
      Po oddaniu 20\,L — tytuł MZ „Honorowy Dawca Krwi — Zasłużony dla Zdrowia Narodu". \\
  \addlinespace
  \term{RCKiK}
    & Regionalne Centrum Krwiodawstwa i Krwiolecznictwa — jednostka nadrzędna
      w danym województwie. \\
  \addlinespace
  \term{Karencja}
    & Ustawowo zdefiniowany, bezwzględnie wymagany czas przerwy między kolejnymi
      donacjami, chroniący zdrowie dawcy. \\
  \addlinespace
  \term{Gamifikacja}
    & Wykorzystanie mechanik znanych z gier (punkty, odznaki, paski postępu, rankingi)
      w celu budowania nawyków i zwiększania zaangażowania użytkowników. \\
  \addlinespace
  \term{OCR}
    & \textit{Optical Character Recognition} — algorytmy przekształcające obraz (zdjęcie,
      skan) w edytowalny ciąg tekstowy; kluczowe do odczytywania wyników morfologii. \\
  \bottomrule
\end{tabularx}
\end{center}

% ════════════════════════════════════════════════════════════
\section{Analiza Interesariuszy}
% ════════════════════════════════════════════════════════════

\begin{center}
\begin{tabularx}{\linewidth}{L{2.6cm} X X L{3cm}}
  \toprule
  \textbf{Kategoria} & \textbf{Opis i Charakterystyka} & \textbf{Główne Oczekiwania} & \textbf{Potencjalne Przyczyny Niezadowolenia} \\
  \midrule
  Dawcy Kluczowi (Regularni ZHDK)
    & Osoby oddające krew wielokrotnie, silnie zmotywowane, ukierunkowane
      na korzyści (zniżki w komunikacji, darmowe leki, obsługa bez kolejki).
    & Bezbłędny licznik sumujący całkowity litraż. Powiadomienia ułatwiające
      zdobycie kolejnych odznak. Precyzyjne generowanie podsumowań do PIT/O.
    & Błędy w sumowaniu litrów wynikające z braku uwzględnienia przeliczników
      dla różnych składników krwi. \\
  \addlinespace
  Dawcy Sporadyczni / Nowi
    & Młode osoby poszukujące cyfrowych, natychmiastowych wzmocnień
      motywacyjnych i informacji edukacyjnych.
    & Nowoczesny interfejs UI, mechaniki grywalizacyjne (zdobywanie poziomów
      i wirtualnych kropli), mapy pokazujące centra poboru.
    & Niezrozumiały, przeładowany interfejs kliniczny. \\
  \addlinespace
  Kobiety — Dawczynie
    & Wysoce świadoma grupa, narażona na odrzucenia lekarskie ze względu
      na wahania poziomu żelaza związane z fizjologią.
    & Inteligentna, dyskretna integracja kalendarza cyklu miesiączkowego,
      wstrzymująca alerty i przesuwająca datę możliwej donacji na okno
      pełnej regeneracji organizmu.
    & Brak mechanizmów różnicujących płeć w kalkulatorach karencji
      i stopniach honorowego krwiodawcy. \\
  \addlinespace
  Stacje Krwiodawstwa (RCKiK)
    & Podmioty zarządzające publiczną służbą krwi, audytorzy poprawności
      stanów magazynowych.
    & Narzędzie, które poprzez precyzyjne powiadomienia zmniejszy wahania
      i deficyty magazynowe grup rzadkich (np. AB Rh$-$).
    & Aplikacja prezentująca dane lub porady medyczne sprzeczne
      z aktualnymi wytycznymi Ministerstwa Zdrowia. \\
  \bottomrule
\end{tabularx}
\end{center}

% ════════════════════════════════════════════════════════════
\section{Analiza Konkurencji i Luka Rynkowa}
% ════════════════════════════════════════════════════════════

\subsection{Rynek Polski: ograniczenia funkcjonalne i architektoniczne}

\subsubsection*{1. Portal Dawcy „Moja Krew" (RCKiK Radom i systemy ogólnopolskie)}
\begin{description}[leftmargin=2.8cm, style=sameline]
  \item[Mocne strony:] Aplikacja połączona bezpośrednio z oficjalnymi bazami danych centrów
    (e-krew). Posiada podgląd stanów magazynowych krwi oraz cenną funkcję sprawdzania
    aktualnego stanu kolejek w punktach stacjonarnych.
  \item[Słabe strony:] Narzędzie określane jako przeładowane informacjami, mało intuicyjne,
    o zachowawczym projekcie graficznym. Użytkownicy wskazują na braki w angażowaniu ich do
    powrotu. Brak jest rozbudowanych modułów cyfrowego archiwum zdrowotnego.
\end{description}

\subsubsection*{2. Aplikacja „Bliscy Krewni" (RCKiK Katowice)}
\begin{description}[leftmargin=2.8cm, style=sameline]
  \item[Mocne strony:] Dobrze zorganizowany moduł społeczny i lojalnościowy
    (,,Lokale Bliskich Krewnych''), pozwalający wylegitymować się aplikacją celem uzyskania
    realnych, stacjonarnych zniżek w kawiarniach i instytucjach kultury.
  \item[Słabe strony:] Silnie regionalny charakter — użyteczność aplikacji spada drastycznie
    poza obszarem Śląska. Doświadcza problemów stabilnościowych, przez co jest nisko
    oceniana w sklepach z aplikacjami.
\end{description}

\subsubsection*{3. Prywatne aplikacje (np. „Dawca Krwi" iOS)}
\begin{description}[leftmargin=2.8cm, style=sameline]
  \item[Słabe strony:] Opierają się na manualnym wpisywaniu informacji. Często monetyzowane
    są poprzez reklamy wewnątrz aplikacji.
\end{description}

% ────────────────────────────────────────────────────────────
\subsection{Wnioski i identyfikacja luk}

Żadna z przebadanych aplikacji nie posiada kompleksowego, łatwego w obsłudze cyfrowego
archiwum badań morfologicznych działającego w oparciu o silnik sztucznej inteligencji (OCR).
Użytkownicy zmuszeni są zbierać papierowe świstki, aby monitorować długofalowe trendy
poziomu żelaza.

Ponadto, na rynku istnieje kompletna próżnia w zakresie uwzględniania biologii
i endokrynologii kobiet — żadne oprogramowanie krwiodawcze nie integruje się z
kalendarzem cyklu menstruacyjnego.

% ════════════════════════════════════════════════════════════
\section{Architektura Funkcjonalności i Priorytety (MoSCoW)}
% ════════════════════════════════════════════════════════════

\subsection{Jądro Aplikacji (\textit{Must-Have})}

\begin{itemize}
  \item \textbf{Multiwymiarowe odliczanie do donacji:} Algorytmy obsługują różne przerwy —
        8~tygodni dla krwi pełnej, 4 lub 2~tygodnie dla płytek krwi oraz osocza,
        krzyżując interwały w zależności od donacji mieszanych i weryfikując z płcią
        użytkownika.

  \item \textbf{Interfejs Profilowy (Karta Dawcy):} Ekran główny eksponujący grupę krwi
        oraz cyfrowy, centralnie ułożony, rosnący licznik oddanych litrów (główny
        wskaźnik sukcesu).

  \item \textbf{Nawigacja:} Zintegrowana mapa prezentująca Centra Krwiodawstwa (RCKiK).
\end{itemize}

\subsection{Rzeczy warte wprowadzenia (\textit{Should-Have})}

\begin{itemize}
  \item \textbf{Alert Zapotrzebowania Grupowego:} Odbieranie powiadomień Push opartych
        o stan magazynów w danym województwie, gdy wystąpi krytyczny niedobór
        zadeklarowanej przez użytkownika grupy krwi.

  \item \textbf{Moduł Podatkowy (Odliczenie PIT):} Algorytm kalkulatora finansowego
        przeliczający zsumowane donacje z danego roku na realną stawkę ulgi podatkowej
        obowiązującą w latach 2025/2026 (130 zł za litr), pomagając dawcy w
        uzupełnianiu rocznego załącznika PIT/O.

  \item \textbf{Ścieżka Progresji ZHDK:} Odliczanie litrów potrzebnych do zdobycia
        kolejnych odznak PCK i Ministerstwa Zdrowia (np. brązowa odznaka 5\,L/6\,L,
        złota 15\,L/18\,L). Aplikacja zaprezentuje na poszczególnych ,,poziomach''
        informacje o odblokowanych przywilejach (zniżki PKP, leki ZK, darmowe
        przejazdy miejskie).
\end{itemize}

\subsection{Innowacje Opcjonalne (\textit{Could-Have})}

\begin{itemize}
  \item \textbf{Cyfrowe Archiwum Morfologiczne z OCR:} Wykorzystanie aparatu smartfona
        i sieci neuronowych (biblioteki rozpoznawania obrazu lub wbudowane silniki
        iOS/Android) do wykonania skanu wydruku badania laboratoryjnego. Aplikacja
        automatycznie rozpozna i wyodrębni wskaźniki (np. Hemoglobina, Krwinki białe,
        Żelazo), zapisując je w bazie i generując wykres historycznych wahań.

  \item \textbf{Funkcjonalność dodawania znajomych:} Możliwość porównania ilości
        oddanej krwi ze znajomymi i wprowadzenie elementu rywalizacji.
\end{itemize}

% ════════════════════════════════════════════════════════════
\section{Specyfikacja Wymagań Funkcjonalnych (Wzorzec Volere)}
% ════════════════════════════════════════════════════════════

Poniżej przedstawiono karty wymagań sporządzone zgodnie z notacją Volere.

% ── Requirement Card W_FUN_01 ────────────────────────────────
\begin{mdframed}[style=reqheader]
  \textbf{Karta:\ \req{W\_FUN\_01}} \hfill
  \textit{Obliczenia Karencji Donacyjnych}
\end{mdframed}
\begin{mdframed}[style=reqcard]
  \begin{tabbing}
    \hspace{3.2cm}\=\kill
    \textbf{Typ}\> Jądro Aplikacji (Wymaganie Funkcjonalne) \\[4pt]
    \textbf{Opis}\> System musi obliczyć i aktywnie odliczać dni do momentu, w~którym \\
                \> dawca prawnie i medycznie może powrócić do RCKiK. \\[4pt]
    \textbf{Przesłanka}\> Okresy są silnie uwarunkowane prawnie. Zależnie od zadeklarowanej \\
                       \> płci oraz ostatniego składnika donacji stosuje się inną formułę \\
                       \> upływu czasu (od 48~godzin do 8~tygodni). \\[4pt]
    \textbf{Kryterium}\> Tester, definiując profil jako ,,Kobieta'', rejestruje donację krwi \\
    \textbf{spełnienia}\> pełnej w dniu 01.03.2026. System musi wyświetlić datę \\
                       \> odblokowania na \textbf{26.04.2026} (dokładnie 8~tygodni przerwy) \\
                       \> i uniemożliwić jej wprowadzenie wcześniejszej donacji krwi pełnej. \\
                       \> Licznik musi zablokować donacje krwi pełnej na 12~miesięcy po \\
                       \> odnotowaniu 4.~donacji w danym roku kalendarzowym.
  \end{tabbing}
\end{mdframed}

% ── Requirement Card W_FUN_02 ────────────────────────────────
\begin{mdframed}[style=reqheader]
  \textbf{Karta:\ \req{W\_FUN\_02}} \hfill
  \textit{Integracja z Cyklem Menstruacyjnym i Blokady Anemiczne}
\end{mdframed}
\begin{mdframed}[style=reqcard]
  \begin{tabbing}
    \hspace{3.2cm}\=\kill
    \textbf{Typ}\> Opcjonalne (Wymaganie Funkcjonalne / Architektura Zdrowotna) \\[4pt]
    \textbf{Opis}\> Kobiety-dawczynie muszą posiadać mechanizm wprowadzania informacji \\
                \> o wystąpieniu i czasie trwania cyklu miesiączkowego, celem ingerencji \\
                \> w globalny kalendarz odliczający karencję. \\[4pt]
    \textbf{Przesłanka}\> Oddawanie krwi podczas lub krótko po krwawieniu grozi ostrą \\
                       \> niedokrwistością i skutkuje dyskwalifikacją przez personel lekarski. \\
                       \> System bez tej funkcjonalności wysyłałby zgubne powiadomienia, \\
                       \> zachęcając do donacji w najgorszym fizjologicznym momencie. \\[4pt]
    \textbf{Kryterium}\> Użytkowniczka zaznacza okres menstruacji na dni 10.05.2026–15.05.2026. \\
    \textbf{spełnienia}\> Jeśli zwykłe wyliczenie karencji pozwalało na powrót 12.05.2026, \\
                       \> algorytm automatycznie nakłada barierę, nadpisując termin donacji na \\
                       \> \texttt{[Koniec\_Miesiączki] + 3 dni\_regeneracji}, \\
                       \> odblokowując system dopiero \textbf{19.05.2026}, uciszając jednocześnie \\
                       \> wszelkie notyfikacje Push do tego czasu.
  \end{tabbing}
\end{mdframed}

% ── Requirement Card W_FUN_03 ────────────────────────────────
\begin{mdframed}[style=reqheader]
  \textbf{Karta:\ \req{W\_FUN\_03}} \hfill
  \textit{Obliczanie Ulgi Podatkowej (Moduł PIT)}
\end{mdframed}
\begin{mdframed}[style=reqcard]
  \begin{tabbing}
    \hspace{3.2cm}\=\kill
    \textbf{Typ}\> Warto Wprowadzić (Wymaganie Funkcjonalne / Biznesowe) \\[4pt]
    \textbf{Opis}\> Aplikacja musi podsumować sumę donacji z danego roku podatkowego \\
                \> i przetworzyć je na wartość ulgi wyrażoną w PLN, z uwzględnieniem \\
                \> algorytmicznych limitów. \\[4pt]
    \textbf{Przesłanka}\> Dawcy uzyskujący dochody opodatkowane (PIT-37, PIT-36, PIT-28) \\
                       \> często nie wiedzą jak poprawnie skalkulować ulgę w załączniku PIT/O. \\
                       \> Realny zastrzyk finansowy jest obok satysfakcji najsilniejszym \\
                       \> zjawiskiem gamifikacyjnym i motywacyjnym. \\[4pt]
    \textbf{Kryterium}\> Dla donacji zgromadzonych w roku 2025 i 2026, aplikacja przemnaża \\
    \textbf{spełnienia}\> wpisane przez użytkownika litry (np. 1,8\,L) przez ustawową stawkę \\
                       \> zdefiniowaną jako zmienna systemowa: \textbf{130\,PLN/litr}. \\
                       \> Wyświetlony wynik \textbf{234,00\,PLN} jest opatrzony formułką \\
                       \> przypominającą, że ostateczna odliczona wartość w~deklaracji nie \\
                       \> może przekroczyć 6\% rocznego dochodu/przychodu dawcy.
  \end{tabbing}
\end{mdframed}

% ── Requirement Card W_FUN_04 ────────────────────────────────
\begin{mdframed}[style=reqheader]
  \textbf{Karta:\ \req{W\_FUN\_04}} \hfill
  \textit{Ścieżka Gamifikacji, Odznaki i Przywileje ZHDK}
\end{mdframed}
\begin{mdframed}[style=reqcard]
  \begin{tabbing}
    \hspace{3.2cm}\=\kill
    \textbf{Typ}\> Warto Wprowadzić (Wymaganie Funkcjonalne / Systemy Lojalnościowe) \\[4pt]
    \textbf{Opis}\> Po każdej donacji system musi operować graficznymi paskami postępu \\
                \> wypełniającymi się aż do uzyskania kolejnego ustawowego progu ZHDK. \\[4pt]
    \textbf{Przesłanka}\> Dawcy, znając konkretny benefit powiązany z progiem \\
                       \> (np. 18\,L u mężczyzny na darmową komunikację miejską), \\
                       \> wykazują silną skłonność do skrócenia przerw między wizytami \\
                       \> w RCKiK. \\[4pt]
    \textbf{Kryterium}\> Pasek postępu dla mężczyzny z historią 11,0\,L powinien wskazywać \\
    \textbf{spełnienia}\> stan osiągnięty jako \textbf{Brązowa Odznaka (6\,L)} i wizualizować \\
                       \> brakujący 1,0\,L do \textbf{Srebrnej Odznaki (12\,L)}. \\
                       \> Po naciśnięciu wskaźnika poziomu ,,Srebro'', na ekranie powinna \\
                       \> ukazać się lista przywilejów (obsługa bez kolejki w aptekach, \\
                       \> darmowe leki dla ZHDK według uprawnień ,,ZK'').
  \end{tabbing}
\end{mdframed}

% ── Requirement Card W_FUN_05 ────────────────────────────────
\begin{mdframed}[style=reqheader]
  \textbf{Karta:\ \req{W\_FUN\_05}} \hfill
  \textit{OCR i Cyfrowe Archiwum Morfologiczne}
\end{mdframed}
\begin{mdframed}[style=reqcard]
  \begin{tabbing}
    \hspace{3.2cm}\=\kill
    \textbf{Typ}\> Opcjonalne (Wymaganie Funkcjonalne / Integracja Modułów Zewnętrznych AI) \\[4pt]
    \textbf{Opis}\> Użytkownik wykorzystuje aparat systemowy z poziomu aplikacji \\
                \> w celu skanowania obrazu fizycznego wydruku badań laboratoryjnych \\
                \> uzyskanych w RCKiK przed donacją, w celu zasilenia bazy trendów \\
                \> zdrowotnych. \\[4pt]
    \textbf{Przesłanka}\> OCR w medycynie eliminuje konieczność manualnego wpisywania \\
                       \> dziesiątek wskaźników przez pacjenta, minimalizując błędy ludzkie. \\
                       \> Krwiodawcy traktują regularne oddawanie krwi jako formę darmowych, \\
                       \> cyklicznych badań profilaktycznych. \\[4pt]
    \textbf{Kryterium}\> Mechanizm OCR parsując zrobione w dobrym oświetleniu zdjęcie formatu \\
    \textbf{spełnienia}\> A4, przy użyciu wyrażeń regularnych (Regex), musi rozpoznać tag \\
                       \> \texttt{HGB} (Hemoglobina) i pomyślnie zlokalizować sąsiadującą \\
                       \> zmiennoprzecinkową wartość cyfrową (np. \texttt{14.2}), po czym \\
                       \> zaproponować użytkownikowi naniesienie jej jako nowego punktu na \\
                       \> osi czasu wykresu morfologicznego w zaledwie 3~kliknięciach.
  \end{tabbing}
\end{mdframed}

% ════════════════════════════════════════════════════════════
\section{Wymagania Niefunkcjonalne i Stos Technologiczny}
% ════════════════════════════════════════════════════════════

\subsection{Wymagania dotyczące używalności i estetyki}

\begin{itemize}
  \item \textbf{Wizualna tożsamość:} Interfejs musi budzić zaufanie, w tonacjach bieli
        i czerwieni (konotacje medyczne, zbliżone do UI Red Cross i Moja Krew).
        Centralny, dominujący na ekranie licznik litrów oraz odliczanie do nowej donacji
        muszą być zrealizowane czytelnymi wskaźnikami.

  \item \textbf{Ergonomia dla seniorów ZHDK:} Z uwagi na starzejącą się pulę
        długoletnich dawców (Zasłużonych dla Zdrowia Narodu, powyżej 20~litrów krwi),
        czcionki muszą ulegać dynamicznemu skalowaniu w odpowiedzi na systemowe
        ustawienia ułatwień dostępu na Android (Accessibility standards).

  \item \textbf{Nawigacja Map:} Przeglądanie akcji lokalnych na mapie nie może być
        blokowane powolnym wczytywaniem bazy danych; należy wdrożyć clusterowanie
        znaczników geolokalizacyjnych stacji krwiodawstwa w przypadku oddalenia
        obrazu (\textit{zoom out}).
\end{itemize}

\subsection{Architektura Wykonawcza i Stos Technologiczny}

\begin{center}
\begin{tabularx}{\linewidth}{L{3.5cm} X}
  \toprule
  \textbf{Warstwa} & \textbf{Technologia i uzasadnienie} \\
  \midrule
  Frontend mobilny
    & \textbf{Flutter} (cross-platform) — generowanie jednoścież kowego kodu
      na platformę Android. Umożliwia implementację UI/UX z natywną
      wydajnością. \\
  \addlinespace
  AI / OCR
    & \textbf{Google ML Kit} (Android) — bezpłatne, działające offline
      procesowanie bloków tekstu z fotografii. Mityguje obciążenie serwera
      chmurowego kosztownymi wyliczeniami obrazu. \\
  \addlinespace
  Backend / Powiadomienia
    & \textbf{Firebase BaaS} — szybkie wdrożenie usług Cloud Messaging
      (wysyłanie powiadomień Push o brakach krwi z centrali) oraz
      logowania użytkowników. \\
  \bottomrule
\end{tabularx}
\end{center}

% ════════════════════════════════════════════════════════════
\section{Podsumowanie Wniosków Projektowych}
% ════════════════════════════════════════════════════════════

W projekcie skupiamy się na sektorze \textit{Digital Health}, ponieważ zauważyliśmy,
że dzisiejszy krwiodawca to nie tylko ,,pacjent'', ale przede wszystkim świadomy
użytkownik smartfona. Oczekuje on konkretów: przejrzystych raportów, informacji
o odznaczeniach i łatwego dostępu do benefitów, takich jak ulgi PIT czy
programy lojalnościowe.

\subsubsection*{Co rozwiązujemy?}

Nasza aplikacja uderza prosto w problemy, z którymi nie radzą sobie przestarzałe,
państwowe portale e-zdrowia (biurokracja i przestarzały interfejs). Kluczowe cechy
systemu:

\begin{itemize}
  \item \textbf{Precyzyjne algorytmy:} Serce systemu to liczniki donacji, które
        bezbłędnie uwzględniają płeć i typ krwi.

  \item \textbf{Innowacyjne funkcje (AI \& OCR):} Zgodnie z analizą MoSCoW,
        wzbogaciliśmy projekt o dodatkowe funkcje — m.in. skanowanie wyników badań
        laboratoryjnych przy użyciu sztucznej inteligencji (OCR).

  \item \textbf{Bezpieczeństwo i etyka:} Wprowadziliśmy inteligentne blokady przed
        zbyt wczesnym oddaniem krwi (np. uwzględniając cykl menstruacyjny u kobiet).

  \item \textbf{Grywalizacja:} Wdrożyliśmy system poziomów, co zamienia oddawanie krwi
        w angażującą ścieżkę użytkownika.
\end{itemize}

% ════════════════════════════════════════════════════════════
\section*{Bibliografia}
% ════════════════════════════════════════════════════════════
\addcontentsline{toc}{section}{Bibliografia}

\begin{enumerate}[label={[\arabic*]}]
  \item \textit{Optimizing the Blood Donation App with Gamification Using User-Centered Design},
        ResearchGate, 2024. \url{https://www.researchgate.net/publication/377518409}

  \item \textit{Dopuszczalne odstępy między donacjami krwi pełnej i jej składników},
        Krwiodawcy.org. \url{https://krwiodawcy.org/dopuszczalne-odstepy-miedzy-donacjami}

  \item \textit{Darowizny na krwiodawstwo PIT}, podatki.gov.pl.
        \url{https://www.podatki.gov.pl/ulgi-i-odliczenia/darowizny-na-krwiodawstwo-pit/}

  \item \textit{Kalendarz donacji}, Krwiodawstwo.pl.
        \url{https://krwiodawstwo.pl/kalendarz-donacji/}

  \item \textit{Program lojalnościowy | Twoja krew, moje życie}, twojakrew.pl.
        \url{https://www.twojakrew.pl/program-lojalnosciowy/}

  \item \textit{Korzyści z HDK — Honorowe krwiodawstwo i krwiolecznictwo}, Krwiodawcy.org.
        \url{https://krwiodawcy.org/korzysci-z-hdk}

  \item \textit{Odznaki Zasłużonego Honorowego Dawcy Krwi}, Krwiodawcy.org.
        \url{https://krwiodawcy.org/odznaki/odznaki-zhdk}
\end{enumerate}

\end{document}